\chapter*{Lampiran}
\addcontentsline{toc}{chapter}{Lampiran}

\section*{A. Rubrik Penilaian Tugas Praktikum (Ringkas)}

\begin{longtable}{|>{\raggedright\arraybackslash}p{3.2cm}|>{\raggedright\arraybackslash}p{4.2cm}|>{\raggedright\arraybackslash}p{6.6cm}|}
\hline
\textbf{Aspek} & \textbf{Bobot indikatif} & \textbf{Indikator} \\
\hline
\endfirsthead
\hline
\textbf{Aspek} & \textbf{Bobot indikatif} & \textbf{Indikator} \\
\hline
\endhead
\hline
\endfoot

Kebenaran program & 40\% & Program berjalan sesuai spesifikasi; tidak ada kesalahan sintaks yang fatal \\
\hline
Struktur kode & 25\% & Pemakaian fungsi/struktur kontrol wajar; tidak ada duplikasi berlebihan \\
\hline
Gaya dan dokumentasi & 15\% & Penamaan jelas; komentar pada bagian non-trivial \\
\hline
Laporan / log & 20\% & Langkah jelas; screenshot atau cuplikan output sesuai permintaan \\
\hline
\end{longtable}

\section*{B. Daftar Referensi dan Sumber Belajar Terbuka}

Sumber berikut mendukung penyusunan modul ini (lihat juga berkas \texttt{sumber\_materi.md} pada repositori modul):

\begin{enumerate}[leftmargin=*]
  \item UNESA. \textit{RPS Pemrograman Dasar} (S1 Sains Data). \url{https://sindig.unesa.ac.id/rps-pdf/s1-sains-data/pemrograman-dasar.pdf}
  \item UB, Departemen Statistika/Sains Data. \textit{RPS Pemrograman Dasar}. \url{https://statistika.ub.ac.id/wp-content/uploads/2024/01/RPS-Pemrograman-Dasar.pdf}
  \item UNS. \textit{OCW --- Algoritme \& Pemrograman Dasar dengan Python}. \url{https://ocw.uns.ac.id/site/detailmakul?params=14b317A4054yaVFVGVU16SXdNek6c23bcbb12f578660dbe4043d851963ceb72c9deE0Zkh4Tk1ERXg%3D}
  \item GitHub: \textit{twseptian/pemrograman-python}. \url{https://github.com/twseptian/pemrograman-python}
  \item Univ. Pancasila. \textit{Modul praktikum Python} (PDF). \url{https://teknik.univpancasila.ac.id/labkom/praktikum/Modul%20Praktikum/Pyhton/modul%20praktikum%20python.pdf}
  \item UIMA Silab. \textit{Modul pembelajaran pemrograman Python}. \url{https://silab.uima.ac.id/wp-content/uploads/2024/10/MODUL-PEMBELAJARAN-PEMROGRAMAN-PHYTON.pdf}
  \item UM. \textit{Modul pengenalan Python} (praktikum, 2016). \url{https://elektro.um.ac.id/wp-content/uploads/2016/04/MODUL-4-PENGENALAN-PYTHON.pdf}
  \item Python Software Foundation. \textit{The Python Tutorial} (EN). \url{https://docs.python.org/3/tutorial/index.html}
  \item Python Software Foundation. \textit{The Python Tutorial} (ID). \url{https://docs.python.org/id/3/tutorial/index.html}
  \item Python Wiki. \textit{BeginnersGuide}. \url{https://wiki.python.org/moin/BeginnersGuide}
  \item Severance, C. \textit{Python for Everybody}. \url{https://www.py4e.com/}
  \item Sweigart, A. \textit{Automate the Boring Stuff with Python}. \url{https://automatetheboringstuff.com/}
  \item Downey, A. B. \textit{Think Python, 2e}. \url{https://greenteapress.com/wp/think-python-2e/}
  \item Pilgrim, M. \textit{Dive Into Python 3}. \url{https://diveintopython3.net/}
  \item Google. \textit{Python Class}. \url{https://developers.google.com/edu/python}
  \item W3Schools. \textit{Python Tutorial}. \url{https://www.w3schools.com/python/}
  \item Programiz. \textit{Python Programming Tutorial}. \url{https://www.programiz.com/python-programming}
  \item LearnPython.org. \url{https://www.learnpython.org/}
  \item Real Python. \url{https://realpython.com/}
\end{enumerate}

\section*{C. Evaluasi Akhir dan Integrasi Kompetensi}

Untuk menilai pencapaian secara menyeluruh, dosen dapat menggabungkan: (1) portofolio tugas praktikum per modul; (2) ujian praktikum tengah dan akhir sesuai RPS \textbf{TIF-105}; (3) mini-proyek akhir yang memadukan struktur kontrol, fungsi, struktur data, akses berkas, serta (opsional) OOP dan tes. Rubrik komprehensif memetakan indikator ke Sub-CPMK pada \texttt{silabus\_praktikum\_python.tex}.

\section*{D. Komponen Penilaian Semester (selaras RPS TIF-105)}

\begin{tabular}{|l|c|p{7cm}|}
\hline
\textbf{Komponen} & \textbf{Bobot} & \textbf{Catatan / pemetaan} \\
\hline
Tugas praktikum mingguan & 28\% & CPMK-1 s.d. CPMK-7 sesuai minggu \\
\hline
Laporan / log praktikum & 17\% & Bukti Sub-CPMK; format sesuai dosen \\
\hline
Kinerja dan rubrik sesi & 10\% & Sikap, K3, partisipasi (CPL-4); observasi lab \\
\hline
Ujian praktikum tengah & 20\% & CPMK-1--4; Minggu 14 \\
\hline
Ujian praktikum akhir / proyek & 25\% & CPMK-5--7; integrasi; Minggu 15--16 \\
\hline
\textbf{Jumlah} & \textbf{100\%} & \\
\hline
\end{tabular}

\section*{E. Lembar Observasi Asisten (contoh)}

Nama mahasiswa / NIM: \ldots\quad Tanggal: \ldots\quad Bab / Sub-CPMK: \ldots

\begin{itemize}[leftmargin=*, label=$\square$]
  \item K3 dan tata tertib laboratorium dipatuhi (CPL-4)
  \item Menyelesaikan prosedur kerja sesuai modul tanpa intervensi berlebihan
  \item Menunjukkan pemahaman konsep (bukan hanya menyalin kode)
  \item Artefak tersimpan dan dinamai sesuai konvensi
\end{itemize}

\noindent Komentar: \hrulefill

\section*{F. Template Penamaan dan Pengumpulan}

\begin{itemize}[leftmargin=*]
  \item Skrip Python: \texttt{NIM\_TIF105\_M\textit{nn}\_tugas.py}
  \item Notebook: \texttt{NIM\_TIF105\_M\textit{nn}.ipynb}
  \item Proyek akhir: folder \texttt{NIM\_TIF105\_proyek/} berisi \texttt{README.md} dan kode sumber;\newline \texttt{requirements.txt} jika ada dependensi.
\end{itemize}

\section*{G. Glosarium Singkat}

\begin{description}[style=nextline, widest=Sub-CPMK:, itemsep=0.25em]
  \item[CPMK:] Capaian pembelajaran mata kuliah.
  \item[CPL:] Capaian pembelajaran lulusan program studi.
  \item[IDE:] Lingkungan pengembangan terintegrasi (editor + alat bantu).
  \item[JSON:] Format pertukaran data berbasis teks berstruktur.
  \item[OBE:] Pembelajaran berbasis capaian (\textit{Outcome-Based Education}).
  \item[OOP:] Pemrograman berorientasi objek.
  \item[pip:] Penginstal paket Python.
  \item[REPL:] \textit{Read--Eval--Print Loop} --- mode interaktif interpreter.
  \item[RPS:] Rencana pembelajaran semester.
  \item[Sub-CPMK:] Indikator operasional yang menurunkan CPMK.
  \item[venv:] Modul lingkungan virtual Python.
  \item[PEP 8:] Panduan gaya kode Python.
\end{description}
